\title{Labbook}
\date{07. - 09. September 2026}
\maketitle

\subsection*{Wetter}
\noindent
\textbf{Datum:} 07/09/26 \\
\textbf{Uhrzeit:} 12:35 Uhr \\

\begin{figure}[h]
     \centering
     \includegraphics[scale=0.4]{img/labbook/wetter_hs1.png}
\end{figure}
\begin{itemize}
    \item Teilweise bedeckter Himmel
    \item Kein Niederschlag
\end{itemize}

\subsection*{5.1.1 Calibration (Systemtemperatur)}
USB 1-5, 9-13 \\
Datendatei-Struktur im Programm (RadioUniversePro):

\begin{table}[htbp]
\centering
\begin{tabular}{cccccccc}
\toprule
rep & dev & sb & on\_tbi & off\_tbi & cal\_tbi & flux & SEFD \\
\midrule
1 & bbc01 & u & 12195,30 & 1001,40 & 2634,20 & 00,00 & 00,00 \\
\bottomrule
\end{tabular}
\end{table}

\begin{align}
T_{b}\cdot \Omega_{S} &= T_{A}\Omega_{A} \\
\Omega_{A} &= \frac{\lambda^2}{A_{e}}, \quad \Omega_{S}=\frac{S_{\nu}}{B_{\nu}}, \quad B_{\nu}=\frac{2k_{B}T_{b}}{\lambda^2} \\
S_{\nu} &= \frac{T_{A}\cdot 2k_{B}}{A_{e}} = f(T_{A},A_{e}) \\
A_{e} &= \pi r^{2}\cdot\eta = 1,5^2\pi\cdot0,5 = 1,125\pi \, \text{m}^2 \\
T_{A} &= \frac{S_{\nu}A_{e}}{2k_{B}}
\end{align}

\begin{quote}
$T_{\text{sys,left}} = 98.76 \pm 22.54 \text{ K}$ \\
$T_{\text{sys,right}} = 107.80 \pm 31.65 \text{ K}$ \\
$T_{\text{sun,left}} = 1237.79 \pm 89.62 \text{ K}$ \\
$T_{\text{sun,right}} = 1302.59 \pm 119.16 \text{ K}$
\end{quote}

Mit $S_{\nu}$ der Sonne (aus SWC):
\begin{itemize}
    \item 76 sfu (solar flux units): $76 \cdot 10^{-11} \, \text{Wm}^{-2}\text{Hz}^{-1} = 76 \cdot 10000 \, \text{Jy} = 760000 \, \text{Jy}$
    \item $\Rightarrow T_A = 972,753 \text{ K}$
    \item Störung wegen RFI
\end{itemize}

\subsection*{5.2 Angular Resolution}
USB 1-5, 9-13 \\
Position der Sonne berechnet durch die Steuerungssoftware:

\begin{table}[htbp]
\centering
\begin{tabular}{cccccccc}
\toprule
Source & RA & DEC & AZ & EL & HA & Type & Note \\
\midrule
Sun & 110410.1 & +055810.9 & +1684733.7 & +42 00 48.8 & -003253.9 & Star & East \\
\bottomrule
\end{tabular}
\end{table}

\begin{itemize}
    \item Im Ordner \texttt{cscan} befindet sich der Screenshot.
    \item Im Ordner \texttt{FITS} befinden sich die Daten (lon vs. lat).
\end{itemize}

\subsection*{5.3 Sensitivity of the Instrument}
Antennentemperatur von VirA (BBC 3,11 zeigt geringsten RFI-Einfluss).

Benötigte Zeit für das Mapping:
\begin{itemize}
    \item 28:35.4s (Software)
    \item 29:27.8s (Eigener Timer)
\end{itemize}

\textit{Hinweis: Screenshot auf Ubuntu und nicht Windows! Files im Ordner \texttt{maps} und in \texttt{FITS} als \texttt{imageproJ01 VirA 01}.}

\subsubsection*{TODO: Theoretical Sensitivity}
\begin{align}
\Delta S_{\nu} &= \frac{2k}{A_{e}} \frac{C_{S}T'_{\text{sys}}}{\sqrt{\Delta \nu \tau}} \quad \text{(spektrale Flussdichte)}
\end{align}
Mit $C_{S} \approx 2$, $T'_{\text{sys,left}} = 98.76 \pm 22.54 \text{ K}$, $T'_{\text{sys,right}} = 107.80 \pm 31.65 \text{ K}$, $A_{e} = 1.125\pi$, $\tau = 1\text{s}$, $\Delta \nu = 61\text{kHz}$:
\begin{align}
\Delta S_{\nu,\text{left}} &= 6,248 \cdot 10^{-24} \frac{\text{W}}{\text{m}^2\text{Hz}} = 624,92 \text{ Jy} \\
\Delta S_{\nu,\text{right}} &= 6,82 \cdot 10^{-24} \frac{\text{W}}{\text{m}^2\text{Hz}} = 682,02 \text{ Jy}
\end{align}

Benötigte Integrationszeit $\tau$ zum Detektieren der ausgewählten Quelle mit dem KRT 3 mit einem Signal-Rausch-Verhältnis $S/N = 10$ in jedem Kanal:
\begin{itemize}
    \item Virgo A: Flussdichte $= 211 \pm 11 \text{ Jy}$, Radiogalaxie
    \item Für 5 Kanäle jede Seite ($\Delta \nu = 5 \cdot 61\text{kHz} = 305\text{kHz}$)
    \item $S/N = 10$
    \item Formel: $\tau = \left( \frac{2k\cdot C_{S}\cdot T'_{\text{sys}}}{\Delta S_{\nu}\cdot A_{E}} \right)^{2} \cdot \frac{\Delta\nu}{S/N}$
    \item $\tau_{l} = 4,12\text{s}$
    \item $\tau_{r} = 5,22\text{s}$
\end{itemize}

\subsubsection*{Determination of the Uncertainty of Radiant Flux Measurements}

\subsubsection*{Preliminary Work}
\begin{align}
S_{\nu} &= \frac{P}{4\pi d^2\sqrt{2\pi}\sigma} = 1263,16 \text{ Jy}
\end{align}
Für Mobile Radio on the Moon (übertragen mit 2W, 900MHz, FWHM: 80kHz, $\sigma = 34\text{kHz}$, empfangen auf der Erde).

\begin{align}
S_{\nu} &\propto \nu^{\alpha}, \quad \alpha = -0,7
\end{align}
Bei $\nu = 1420.4\text{ MHz}$: $S_{\nu,\text{CasA}} = 1420 \pm 70 \text{ Jy}$.
\begin{align}
c &= \frac{1420,4^{-0,7}\text{ MHz}}{1420\text{ Jy}} \\
S_{\nu,\text{Cas A}} &= \frac{900^{-0,7}}{c} \text{ Jy} = 1954,37 \text{ Jy} \\
T_{A} &= \frac{S_{\nu}A_{e}}{2k_{B}} = 1,82 \text{ K} \quad (\text{für } \nu = 1420.4\text{ MHz})
\end{align}

\noindent\textbf{TODO: Fehlerfortpflanzung}

\subsubsection*{Evaluation}
BBC 3,11
\begin{figure}[h]
     \centering
    \includegraphics[scale=0.25]{img/labbook/map vira red.png}
\end{figure}
\begin{figure}[h]
     \centering
    \includegraphics[scale=0.3]{img/labbook/ds9 region data.png}
\end{figure}
\begin{figure}[h]
     \centering
    \includegraphics[scale=0.25]{img/labbook/ds9 region screenshot.png}
\end{figure}

\begin{align}
\text{mean} &= 1455,83 \\
\text{std} &= 212,206 \\
\sigma[K] &= \frac{\text{std}}{\text{mean}}\cdot T'_{\text{sys}} = 15,71 \text{ K} \\
\text{mean}_{\text{scaled}}[K] &= T_{A}(\text{OFF}) = 156938,474 \text{ K} \\
T_{A}(\text{ON}) &= 3067 \cdot 107,8\text{K} = 330622,6 \text{ K}
\end{align}

Antennentemperatur der Quelle:
\begin{align}
T_{s} &= \frac{T_{A}(\text{ON})-T_{A}(\text{OFF})+T'_{\text{sys}}}{T'_{\text{sys}}} = 1612,17 \text{ K}
\end{align}

\begin{itemize}
    \item Rauschen ist viel größer als erwartet.
    \item Vergleich mit theoretischer Empfindlichkeit (600 Jy, ca. 1 K).
    \item Bei uns jetzt 128 Kanäle, 10s Integrationszeit $\rightarrow$ eigentlich sollte es noch besser sein ($\ll 1\text{ K}$).
    \item Bei uns jedoch viel Einfluss von RFI $\rightarrow \sim 15\text{ K}$ Std./Rauschen.
    \item Empfänger-/Verstärkersensibilität ist ein Problem; thermisches Rauschen begrenzt nicht, aber besser als 1 K kann es eigentlich nicht werden.
\end{itemize}

Gesucht: $S/N$ (25 Karten für $S/N = 10$).

\subsection*{5.1.2 Background Radiation}
\subsubsection*{Mond}
Elevation: $30^\circ \, 52' \, 32.0''$ (in Grad, Minuten, Sekunden) \\
Für BBC 1-5, 9-13 (wie zuvor für die Sonne, auch bei schlechter Qualität).
\begin{align}
\text{Span} &= 2(\text{ELV} - 10^\circ) = 2(30^\circ - 10^\circ) = 40^\circ
\end{align}
Da der Azimut bei $268^\circ \, 03' \, 32.7''$ liegt, wählen wir einen etwas kleineren Span: $\text{Span} = 38^\circ$. \\
Zeitpunkt: ca. 14:56 Uhr. \\
Daten in FITS: \texttt{LON-MOON} und \texttt{LAT-MOON} + Screenshot in \texttt{cscan}: \texttt{CSCAN-MOON}.

\subsubsection*{Virgo A}
Zeitpunkt: 15:28 Uhr \\
Elevation: $47^\circ \, 56' \, 18''$, Azimut: $196^\circ$, $\text{Span} = 74$.

\subsection*{Observation of Neutral Hydrogen in the Milky Way}
\subsubsection*{Frequency Setting of the Receiver}
\subsubsection*{Preliminary Work}
\begin{itemize}
    \item Was ist die Frequenzauflösung $\delta\nu$ (Bandbreite eines einzelnen Kanals) und wie lautet die Geschwindigkeitsauflösung $\delta\upsilon$ bei der Beobachtung der H I-Linie?
    \item In welchem Kanal würden Sie die H I-Linie bei $\nu_0 = 1420.4\text{ MHz}$ erwarten, unter Annahme einer Radialgeschwindigkeit von $0\text{ km s}^{-1}$?
    \item Sie wollen das Radiokontinuum einer Quelle messen. Welche Frequenzen sind innerhalb der Bandbreite erlaubt, damit die Messung nicht durch die 21cm-Linie kontaminiert wird? (Maximale Radialgeschwindigkeit des H I: $\pm 200\text{ km s}^{-1}$).
    \item Welche Frequenzbänder sollten in BBC Tools ausgewählt werden, wenn man einmal die H I-Linie und einmal das Kontinuum beobachten möchte?
\end{itemize}

\subsubsection*{Measurement}
\begin{align}
R &= R_{0}\cdot \sin(l), \quad R_{0} = 8,5\text{ kpc}
\end{align}

\begin{table}[htbp]
\centering
\begin{tabular}{cccc}
\toprule
$l$ [deg] & $R$ [kpc] & $v_{\max}$ [km/s] & $v_{R}$ [km/s] \\
\midrule
10 & $1,48 \pm 4,09$ & -- & -- \\
20 & $2,91 \pm 4,0$ & $70,70 \pm 14,50$ & $\dots$ \\
\bottomrule
\end{tabular}
\end{table}

\begin{itemize}
    \item $v_{\max}$ grafisch bestimmt (Schnittwert mit der orangen $\sigma$-Linie).
    \item $v_{R} = v_{\max} + \omega_{0}R_{0}\cdot \sin(l)$ mit $\omega_{0}R_{0} = 220\text{ km/s}$.
    \item $v_{\max,\text{err}}$ grafisch über blauen Bereich bestimmt (Schnittpunkt min -- Schnittpunkt max / 2).
    \item $\Delta l = \pm 2^\circ$.
    \item Fehlerfortpflanzung für $R$:
    \begin{align}
    \Delta R &= \left|\frac{\partial R}{\partial l}\right|\cdot \Delta l = |R_{0}\cdot \cos(l)|\cdot \Delta l \\
    \sigma_{R} &= \sqrt{|R_{0}\cos(l)|\Delta l}
    \end{align}
    \item Eingeschlossene Masse:
    \begin{align}
    M(R) &= 0,234 \cdot 10^{10} \left( \frac{v(R)}{100\text{ km/s}} \right)^2 \left( \frac{R}{\text{kpc}} \right) M_{\odot}, \quad 1 M_{\odot} = 1,98 \cdot 10^{30}\text{ kg}
    \end{align}
\end{itemize}

\subsubsection*{Distribution of Dark Matter}
\begin{align}
M_{\text{DM}}(R_0) &= 66411575447.22 \pm 3593349279.00 \, M_{\odot}
\end{align}